% Generated draft: requires mathematical and source review.
\documentclass[11pt]{article}
\usepackage[margin=25mm]{geometry}
\usepackage{fontspec}
\setmainfont{Latin Modern Roman}
\usepackage{amsmath,amssymb,mathtools,unicode-math}
\setmathfont{Latin Modern Math}
\usepackage{enumitem,needspace,xurl,hyperref}
\setlength{\parindent}{0pt}
\setlength{\parskip}{5pt}
\setlength{\emergencystretch}{3em}
\newcommand{\N}{\mathbb{N}}
\newcommand{\Z}{\mathbb{Z}}
\newcommand{\Q}{\mathbb{Q}}
\newcommand{\R}{\mathbb{R}}
\newcommand{\C}{\mathbb{C}}
\newcommand{\F}{\mathbb{F}}
\newcommand{\A}{\mathbb{A}}
\newcommand{\PP}{\mathbb{P}}
\newcommand{\E}{\mathbb{E}}
\newcommand{\eps}{\varepsilon}
\newcommand{\dd}{\,\mathrm{d}}
\begin{document}
\section*{UnsolvedMath ID: 30003875}
\textbf{Source problem code:} catalogue entry
\subsection{Definitions and mathematical statement}
For a field \(k\), an abelian variety over \(k\) is a geometrically connected, smooth, projective algebraic group. A principal polarization on an abelian variety \(A\) is a polarization inducing an isomorphism \(A\to A^\vee\), where \(A^\vee\) is the dual abelian variety. The notation \(A_g\) in the supplied record denotes the moduli space of principally polarized abelian varieties of dimension \(g\). For a positive integer \(p\), the group scheme of \(p\)-torsion is
\[
A[p]=\ker([p]\colon A\longrightarrow A).
\]
The record denotes a corresponding \(p\)-torsion cover by \(A_{g,p}\to A_g\), but does not specify the torsion data parametrized by this cover.

For clarity, a standard definition of resolvent degree for a dominant generically finite map \(f\colon Y\dashrightarrow X\) of irreducible varieties over a fixed field \(k\) is the least nonnegative integer \(d\) for which there is a tower of dominant generically finite rational maps
\[
X_r\dashrightarrow X_{r-1}\dashrightarrow\cdots\dashrightarrow X_0=X
\]
and a dominant rational map \(X_r\dashrightarrow Y\) over \(X\), such that each step is, birationally, a component of a pullback of a dominant generically finite map \(W_i\dashrightarrow V_i\) with \(\dim V_i\leq d\). All varieties and maps in this definition are over \(k\). This definition explains the notation \(\operatorname{RD}_k(f)\); the supplied record does not specify its base field or whether its moduli objects require a stack interpretation.

\textbf{Supplied conjectural statement.} For every pair \((g,p)\) in the intended parameter range, subject to \((g,p)\neq(2,2)\), the proposed equality is
\[
\operatorname{RD}(A_{g,p}\longrightarrow A_g)
=\dim A_g
=\frac{g(g+1)}{2}.
\]
The record does not explicitly give the range of \(g\), state whether \(p\) must be prime, fix a base field or characteristic, or define the particular torsion cover. Consequently, the displayed formula preserves the supplied target, but these unresolved choices must be specified before it becomes a fully determined mathematical assertion. No assertion for \((g,p)=(2,2)\) is supplied.
\subsection{Short English statement}
The proposed statement says that the indicated torsion cover of the moduli space of principally polarized abelian varieties has resolvent degree equal to the full moduli dimension, except at \((g,p)=(2,2)\). The parameter ranges, base field, and precise torsion cover remain unspecified. The supplied record labels the problem open; this draft does not independently verify that status.
\subsection{Sources}
\begin{itemize}
\item Supplied local UnsolvedMath registry record, permanent ID \(30003875\), exact source problem code \texttt{catalogue entry}, titled ``Resolvent Degree of Abelian-Variety Torsion Covers,'' in the Oberwolfach Reports Open Problems collection. The record gives proposed year 2018 and supplies no named proposer.
\item Supplied source citation: \emph{Algebraische Zahlentheorie} (2019), Oberwolfach Reports, DOI 10.4171/owr/2018/30, pp. 1785--1844. \url{https://doi.org/10.4171/owr/2018/30}. The supplied background describes these reports as documenting mathematical workshops and their open-problem discussions.
\item Supplied literature source: ``Oberwolfach Report 30/2018.'' \url{https://ems.press/content/serial-article-files/46751}.
\item Supplied literature source: Farb--Wolfson, ``resolvent degree.'' \url{https://arxiv.org/abs/1803.04063}.
\item Literature assessment dated exactly \texttt{2026-08-21}: the supplied curation labels the problem open, describes maximal resolvent degree here as a conjectural lower-bound statement, and reports that no general proof was identified. This is supplied dated curation derived from the OWR material, not an independently verified literature assessment; its time-sensitive status and cited claims require independent verification.
\end{itemize}
\end{document}
